%% ma: L10-B14-0078
%% lop: 10
%% chuong: 5
%% bai: 14
%% dang: 10.14.9
%% loai: TN4
%% muc: VD
%% dap_an: "C"
%% nguon: "chuyen-de-cac-so-dac-trung-cua-mau-so-lieu-khong-ghep-nhom-toan-10-chuong-trinh-moi.pdf (GV. Phan Nhật Linh) (P05), Bài 03 Dạng 2 Phần I Câu 25"
%% kien_thuc: "Lớp 10 Bài 13 (số trung bình), Bài 14 (phương sai, độ lệch chuẩn)"
%% trang_thai: chinh-thuc
%% ghi_chu: "Đề gốc ghi 'của bạn An và bạn Bình' nhưng bảng có ba bạn, đã sửa. Tự duyệt theo yêu cầu giáo viên 22/10/2026."
\begin{ex}
Kết quả $5$ lần nhảy xa (đơn vị: mét) của ba bạn An, Bình, Cường cho ở bảng sau:
\begin{center}
\begin{tabular}{|c|c|c|c|c|c|}\hline
An & $2{,}1$ & $2{,}5$ & $2{,}6$ & $2{,}8$ & $2{,}9$\\\hline
Bình & $2{,}7$ & $2{,}1$ & $2{,}3$ & $2{,}8$ & $2{,}6$\\\hline
Cường & $2{,}2$ & $2{,}4$ & $2{,}3$ & $2{,}6$ & $2{,}9$\\\hline
\end{tabular}
\end{center}
Bạn nào có kết quả nhảy xa ổn định nhất?
\choice{An}{Bình}{\True Cường}{Như nhau}
\loigiai{$s_A^2=0{,}0776$; $s_B^2=0{,}068$; $s_C^2=0{,}0616$.

Vì $s_C^2<s_B^2<s_A^2$ nên Cường ổn định nhất. Chọn C.}
\end{ex}
